\section{Results}
\label{sec:results}

\subsection{Cohort and evaluation setting}

Table~\ref{tab:cohort} summarizes the cohort and its partition. The
patient-held-out test set comprises 4,107 images from 240 patients across four
centers. The split is performed at the patient level and contains no patient
overlap with the training set; all four centers represented in the test set
also appear in the training split, so the test set characterizes performance
on unseen patients rather than on unseen institutions.

The test cohort contains 60 patients in each of F0, F1, F2 and F3, whereas
training and validation follow the cohort distribution, in which F2 is the most
frequent patient-level grade. Because each of the four stages contributes 60 of
the 240 test patients, aggregate performance should not be interpreted as a
prevalence-weighted estimate for a population screening setting. Image-level
counts are not exactly balanced, as patients contribute differing numbers of
images.

All results below use each configuration's validation-selected checkpoint
under the single prespecified seed~2026.

\subsection{Grading performance}

Table~\ref{tab:main} reports the composite endpoint and its constituent
image-level and patient-level metrics for all five configurations. The observed
ordering by $R_{\mathrm{final}}$ is
E~($0.186865$) $<$ A~($0.187312$) $<$ D~($0.196686$) $<$ B~($0.198100$) $<$
C~($0.200202$). SynAP-Fib (E) therefore attains the best observed
$R_{\mathrm{final}}$ among the five configurations.

Configuration E also records the lowest image-level COR ($0.177908$) and mean
absolute error ($0.377857$), the lowest patient-max COR ($0.200528$) and the
lowest patient-max MAE ($0.420558$). Metric-specific optima differed across
configurations: B yielded the highest four-class accuracy ($67.40\%$), whereas
C yielded the highest $\pm0.5$ accuracy ($70.17\%$); E achieved the lowest
observed $R_{\mathrm{final}}$.

\subsection{Position and lesion priors in combination}

Table~\ref{tab:priors} arranges the four stretch-resize configurations as a
$2\times2$ comparison of the position and lesion priors. Introducing either
prior on its own degraded the composite grading metric relative to the
baseline: $R_{\mathrm{final}}$ rose from $0.187312$ for A to $0.200202$ for the
position-only configuration C and to $0.196686$ for the lesion-only
configuration D. The same direction holds for both constituent COR terms, so
the degradation is not confined to a single aggregation level.

The joint configuration E did not inherit this degradation. It recovered the
loss observed under either single prior and attained the best observed
$R_{\mathrm{final}}$, image COR and patient-max COR of the four
configurations. The advantage of the joint configuration over the position-only
configuration is supported by the supplementary paired bootstrap analysis
(95\% CI for E~$-$~C: $-0.024376$ to $-0.002223$). This is a non-additive
empirical pattern: the effect of supplying both priors together is not what the
two single-prior configurations would suggest in isolation.

\subsection{Stage-wise performance}

Fig.~\ref{fig:stagewise} compares the baseline and SynAP-Fib across the four
clinical stages. Stage-wise MAE was recomputed post hoc from the frozen
per-image predictions.

For stage-wise MAE, E is lower than A at F0 ($0.196277$ versus $0.196508$),
F1 ($0.411035$ versus $0.416107$) and F3 ($0.401280$ versus $0.403570$), and
higher than A at F2 ($0.516525$ versus $0.514474$). For both configurations
the largest stage-wise MAE occurs at F2.

For stage-wise recall, E is higher than A at F0 ($81.14\%$ versus $80.18\%$)
and F1 ($50.50\%$ versus $49.50\%$), equal at F2 ($63.63\%$ for both) and lower
at F3 ($72.10\%$ versus $73.27\%$). For both configurations the lowest recall
occurs at F1.

The two metrics therefore identify different stages as most difficult: F2
carries the largest stage-wise MAE, while F1 carries the lowest recall.

\subsection{Anatomical prior outputs and computational cost}

Beyond the fibrosis score, the joint configuration emits two additional
outputs, summarized in Table~\ref{tab:auxiliary}. Position classification over
the six-class anatomical protocol reaches $66.81\%$ accuracy and $0.6659$
macro-F1 in E, compared with $66.40\%$ and $0.6623$ in the position-only
configuration C, with a lower cross-entropy ($0.9165$ versus $0.9444$).
Agreement with the weak bounding boxes reaches Dice@0.5 of $0.3603$ and
IoU@0.5 of $0.2355$ over 4,035 valid boxes in E, compared with $0.3539$ and
$0.2261$ in the lesion-only configuration D. These lesion figures quantify
agreement with weak boxes and are not segmentation ground-truth performance.

The cost of these additional outputs is modest. Relative to the baseline,
E adds 1,442,673 parameters, a $5.9\%$ increase, and $2.8\%$ more
multiply--accumulate operations, with median single-image latency rising from
$3.031$~ms to $3.993$~ms. Learned gate distributions for both residual pathways
are reported in the Supplementary Material.

\subsection{Checkpoint selection diagnostic}

For configurations A, B, D and E the validation-selected checkpoint also
outperformed that configuration's epoch-120 checkpoint on the test set.
Configuration C is the sole exception, with a best-checkpoint test
$R_{\mathrm{final}}$ of $0.200202$ against $0.191930$ at epoch~120. The
prespecified selection rule was applied without post-hoc substitution, and the
reported ranking therefore uses the validation-selected checkpoint throughout.

Paired bootstrap analyses of all prespecified contrasts, together with
discrimination, calibration, patient-median and per-center results, are
provided in the Supplementary Material.
